\documentclass[a4paper,11pt]{ltjsarticle}

\usepackage[margin=24mm]{geometry}
\usepackage{amsmath,amssymb,mathtools}
\usepackage{booktabs}
\usepackage{longtable}
\usepackage{url}
\usepackage{xcolor}
\usepackage[colorlinks=true,linkcolor=blue!55!black,urlcolor=blue!55!black]{hyperref}

\newcommand{\R}{\mathbb{R}}
\newcommand{\code}[1]{\texttt{#1}}
\newcommand{\RelErr}{\varepsilon}

\title{AlgRemez のコード読解\\
       \large C++ 実装が解いている問題、アルゴリズム、および Julia 移植時の論点}
\author{コード調査メモ}
\date{\today}

\begin{document}
\maketitle

\begin{abstract}
本稿は、\path{upstream AlgRemez/src} にある実装をコードに即して説明するものである。
このプログラムは、正の区間上の
\(f(x)=x^{p/q}\)（および内部 API では指数因子を掛けた関数）に対し、
指定した分子次数と分母次数をもつ有理関数を構成する。
最小化対象は絶対誤差ではなく相対誤差の一様ノルムである。
実装上は、誤差の零点で補間問題を解き、その零点を反復的に移動して、
隣接区間にある誤差極値の絶対値を揃える Remez 型の交換法になっている。
収束後には実根を Newton 法で求め、次数が等しい場合に有理関数とその逆数の
部分分数展開を出力する。
なお、ファイルの拡張子と実装内容から分かるとおり、これは C ではなく C++ のコードである。
\end{abstract}

\tableofcontents

\section{調査対象と全体像}

調査対象は次のファイルである。

\begin{center}
\begin{tabular}{@{}ll@{}}
\toprule
ファイル & 役割 \\
\midrule
\path{upstream/src/main.C} & コマンドライン入力、計算実行、結果のファイル出力 \\
\path{upstream/src/alg_remez.h} & \code{AlgRemez} クラスの宣言 \\
\path{upstream/src/alg_remez.C} & Remez 反復、線形方程式、根、部分分数展開 \\
\path{upstream/src/bigfloat.h} & GMP/MPFR を包む多倍長浮動小数点クラス \\
\path{upstream/src/Makefile} & GMP/MPFR とリンクする古い形式のビルド規則 \\
\path{README} & 利用方法、制限、引用情報 \\
\bottomrule
\end{tabular}
\end{center}

処理の流れは以下である。

\begin{enumerate}
  \item 区間、次数、指数、有効桁数を読み込む。
  \item 誤差零点と誤差極値位置の初期値を Chebyshev 点に似た配置で作る。
  \item 現在の誤差零点で \(R(x_i)=f(x_i)\) となる有理関数を線形方程式から求める。
  \item 各零点間で相対誤差の極値を探索する。
  \item 極値の大きさが等しくなる方向へ零点を移動し、収束まで 3--4 を繰り返す。
  \item 分子・分母の実根を Newton 法と逐次デフレーションで求める。
  \item \(n=d\) の場合、有理関数とその逆数を部分分数展開する。
  \item \path{approx.dat} と \path{error.dat} を出力する。
\end{enumerate}

\section{近似している数学的問題}

\subsection{対象関数}

通常の公開インターフェースが対象とする関数は
\begin{equation}
  f(x)=x^\alpha,\qquad
  \alpha=\frac{p}{q},\qquad p,q\in\mathbb{Z}_{>0},
  \label{eq:target}
\end{equation}
であり、近似区間は
\begin{equation}
  x\in[a,b]
\end{equation}
である。典型用途では \(0<a<b\) を仮定しているが、コード自身はこの条件を検査しない。
負の指数 \(x^{-p/q}\) は直接生成せず、\(x^{p/q}\) の近似の逆数を使う。

引数の多い \code{generateApprox} には、さらに
\begin{equation}
  f(x)=x^{p/q}\exp\!\left(\sum_{j=0}^{L-1}a_jx^{k_j}\right),
  \qquad L\leq 10,
  \label{eq:exp-target}
\end{equation}
を扱う内部機能がある。ただし \path{main.C} は常に \(L=0\) とするため、
コマンドラインからこの機能は使われない。

\subsection{有理関数の表現}

分子次数 \(n\)、分母次数 \(d\) の有理関数
\begin{equation}
 R_{n,d}(x)=\frac{P_n(x)}{Q_d(x)}
 =\frac{\displaystyle\sum_{j=0}^{n}p_jx^j}
        {\displaystyle x^d+\sum_{j=0}^{d-1}q_jx^j}
 \label{eq:rational}
\end{equation}
を使う。分母の最高次係数を 1 に固定して全体のスケールの不定性を除いている。
未知数は
\begin{equation}
 (p_0,\ldots,p_n,q_0,\ldots,q_{d-1})
\end{equation}
の \(K=n+d+1\) 個であり、コードでは \code{param} にこの順番で格納される。
\code{approx} は分子・分母を Horner 法で評価する。

\subsection{最小化する誤差}

\code{getErr} が計算する符号付き相対誤差は
\begin{equation}
  \RelErr(x)=\frac{R_{n,d}(x)-f(x)}{f(x)}
             =\frac{R_{n,d}(x)}{f(x)}-1
  \label{eq:relative-error}
\end{equation}
である。探索ルーチンが実際に返す値は \(E(x)=|\RelErr(x)|\) である。
したがって目標は
\begin{equation}
  \min_{R_{n,d}}\max_{x\in[a,b]}|\RelErr(x)|
  \label{eq:minimax}
\end{equation}
という相対誤差の minimax 近似である。
これは重み \(w(x)=1/f(x)\) をもつ重み付き一様近似と見なせる。

非退化な最良近似では、通常 \(K+1=n+d+2\) 個の点で符号が交互に変わる
等振幅誤差（equioscillation）が期待される。このコードはその状態を、
\emph{誤差の \(K\) 個の零点}を動かすことで作ろうとしている。

\section{Remez 反復の実装}

\subsection{誤差零点での補間}

現在の誤差零点候補を
\begin{equation}
  a<\xi_0<\xi_1<\cdots<\xi_{K-1}<b
\end{equation}
とする。\code{equations} は各点で
\begin{equation}
  R_{n,d}(\xi_i)=f(\xi_i),\qquad i=0,\ldots,K-1
  \label{eq:interpolation}
\end{equation}
を課す。式~\eqref{eq:rational} の分母を払うと、各行は
\begin{equation}
 \sum_{j=0}^{n}p_j\xi_i^j
 - f(\xi_i)\sum_{j=0}^{d-1}q_j\xi_i^j
 = f(\xi_i)\xi_i^d
 \label{eq:linear-row}
\end{equation}
となる。これは未知係数について線形であり、\(K\) 本の連立一次方程式を構成する。

重要なのは、典型的な Remez の説明で用いられる誤差振幅 \(\Delta\) を
未知数として直接解いてはいないことである。現在の零点で対象関数を正確に補間し、
その後に零点を動かすという実装である。

\subsection{初期点}

\code{initialGuess} は Chebyshev 点を基にしつつ、低い \(x\) 側へ点を寄せる変換
\begin{equation}
  \phi(t)=\frac{e^t-1}{e-1},\qquad 0\leq t\leq 1
  \label{eq:warp}
\end{equation}
を使う。初期の極値候補 \(m_i\) は
\begin{align}
 t_i^{(m)} &= \frac{1}{2}\left(1-\cos\frac{\pi i}{K}\right),
       &&i=0,\ldots,K,\\
 m_i &= a+(b-a)\phi\!\left(t_i^{(m)}\right)
 \end{align}
である。初期の零点候補は
\begin{align}
 t_i^{(\xi)} &= \frac{1}{2}\left(1-
       \cos\frac{\pi(2i+1)}{2K}\right),\\
 \xi_i &= a+(b-a)\phi\!\left(t_i^{(\xi)}\right),
       &&i=0,\ldots,K-1
 \end{align}
である。実装は \(i=K\) の値も配列に書くが、反復で使用する零点は
\(\xi_0,\ldots,\xi_{K-1}\) だけである。

指数関数による変換そのものは C++ の \code{double} で計算されてから
多倍長型へ代入される。このため、初期点の精度は指定した多倍長精度ではなく、
実質的に倍精度へ制限される。ただし初期値なので、収束すれば最終精度を
直接制限するとは限らない。

\subsection{各区間の誤差極値探索}

零点によって区間を
\begin{equation}
 [a,\xi_0],\ [\xi_0,\xi_1],\ldots,
 [\xi_{K-2},\xi_{K-1}],\ [\xi_{K-1},b]
\end{equation}
の \(K+1\) 個に分ける。\code{search} は各区間で、前回の極値位置 \(m_i\) と
ステップ \(s_i\) から出発し、\(E(x)=|\RelErr(x)|\) が増える方向へ歩く。
増加しなくなるか、10 ステップ進むか、区間境界に達すると停止する。

これは微分を使った精密な最大化や bracket 付き最適化ではなく、固定ステップの
単純な山登りである。\code{getErr} は誤差の符号も返すが、\code{search} は
その符号を零点更新にも収束判定にも使用しない。したがって、隣接する極値の符号が
実際に交互であることは明示的に検証されない。補間零点が区間内に保たれ、誤差が連続で
各零点を横切るなら交互性が期待できる、という前提に依存している。

\subsection{零点の交換・更新}

探索で得た極値の絶対値を
\begin{equation}
  E_i=|\RelErr(m_i)|,\qquad i=0,\ldots,K
\end{equation}
とする。コードが収束尺度として用いる spread は
\begin{equation}
  S=\frac{\max_i E_i-\min_i E_i}{\min_i E_i}
  \label{eq:spread}
\end{equation}
である。\(\min_i E_i=0\) の場合だけ除算を避ける。

隣接する極値の不均衡から零点移動量を
\begin{equation}
 s_i \leftarrow \eta\,
 \min\!\left(\frac{E_i}{E_{i+1}}-1,\,0.25\right)
 (m_{i+1}-m_i),
 \qquad i=0,\ldots,K-1
 \label{eq:step-update}
\end{equation}
とし、
\begin{equation}
  \xi_i\leftarrow\xi_i-s_i
  \label{eq:zero-update}
\end{equation}
で更新する。ここで緩和係数 \(\eta\) はコード中の変数 \code{delta} であり、
初期値は 0.25 である。これは minimax 誤差振幅を表す記号ではない。
新しい spread が直前より改善しなければ \(\eta\) を半分にする。
更新後の零点が \([m_i,m_{i+1}]\) から出そうな場合は、旧零点と対応する極値との
中点へ戻し、点の順序を保つ。

上側の比には 0.25 のクリップがある一方、コードには対称な下側クリップはない。
ただし \(E_i/E_{i+1}\geq0\) なので、通常は括弧内の下限が \(-1\) となる。

\subsection{停止条件}

初期値 \code{spread=1e37} から反復し、
\begin{equation}
  S\leq 10^{-15}
  \label{eq:stop}
\end{equation}
で停止する。この閾値は要求した多倍長精度とは独立に固定されている。
外側反復には最大回数がない。また、緩和係数 \(\eta\) が \(10^{-15}\) 未満に
なると、精度を増やすよう表示してプロセスを終了する。

返り値の ``error'' は、収束時の全点の最大値を再計算したものではなく、
最初の極値位置 \(m_0\) における \(|\RelErr(m_0)|\) である。
等振幅状態が正しく得られていれば、これは一様誤差振幅と一致する。

\subsection{擬似コード}

コードの本質を簡略化すると次のようになる。

\begin{verbatim}
K = n + d + 1
initialize K error zeros xi[0:K-1]
initialize K+1 extrema guesses m[0:K]
initialize search steps; eta = 0.25
spread = large

while spread > 1e-15
    solve R(xi[i]) = f(xi[i]) for coefficients of P and monic Q

    for each interval separated by xi
        walk from m[i] to find a local maximum of abs(R/f - 1)
        store its location m[i] and magnitude E[i]
    end

    new_spread = (maximum(E) - minimum(E)) / minimum(E)
    halve eta if new_spread did not improve
    move each xi[i] according to E[i] / E[i+1]
    keep xi[i] between m[i] and m[i+1]
end
\end{verbatim}

\section{連立一次方程式の解法}

\code{simq} は式~\eqref{eq:linear-row} の密行列を、行スケーリングを考慮した
部分 pivot 選択付き Gaussian 消去で解く。行 \(i\) の最大絶対値を \(r_i\) とし、
pivot 候補では概ね
\begin{equation}
  \frac{|A_{ik}|}{r_i}
\end{equation}
が最大の行を選ぶ。行自体を物理的に交換せず、\code{IPS} という行番号配列で
置換を表現し、その後に前進・後退代入する。

一方、行列は \(1,x,x^2,\ldots\) という単項式基底で作られる。
区間の幅、次数、スペクトル比 \(b/a\) が大きいと強く悪条件化しやすい。
多倍長演算は丸め誤差を軽減するが、基底そのものの条件数は改善しない。
Julia 版で高次数を重視するなら、区間を \([-1,1]\) に写して Chebyshev 基底を
使う、あるいは列スケーリングを導入する余地がある。

\section{根と部分分数展開}

\subsection{根の計算}

収束後、\code{root} は分子と分母を別々に処理する。多項式 \(H(x)\) の根を
Newton 法
\begin{equation}
  x_{k+1}=x_k-\frac{H(x_k)}{H'(x_k)}
\end{equation}
で一つ求めるごとに一次因子を除き、残った多項式へ同じ処理を繰り返す。
多項式と導関数は Horner 法で評価される。

実装上の重要な制約は次のとおりである。

\begin{itemize}
  \item 実数演算だけであり、複素根は扱えない。
  \item Newton 法へ渡す \([-10^5,1]\) は実際には bracket として強制されず、
        初期値 \((-10^5+1)/2\) を与えるためだけに使われる。
  \item 反復上限は各根 10000 回、停止条件は \(|\Delta x|<10^{-20}\) に固定される。
  \item 根が 0 と返ると失敗扱いなので、真の零根も表現できない。
  \item 逐次デフレーションは、求めた根の誤差を後続の根へ蓄積しうる。
\end{itemize}

根を \(r_j\)、分母の根を \(\rho_j\) とすると、分母が monic なので
\begin{equation}
  R_{n,d}(x)=p_n
  \frac{\prod_{j=1}^{n}(x-r_j)}
       {\prod_{j=1}^{d}(x-\rho_j)}.
  \label{eq:factor}
\end{equation}
コード中の \code{norm} は \(p_n\) である。

\subsection{部分分数展開}

\code{getPFE} と \code{getIPFE} は \(n=d\) の場合だけを扱う。
README および出力が採用する規約は
\begin{equation}
 R_{n,n}(x)=c+\sum_{i=1}^{n}\frac{A_i}{x+\beta_i},
 \qquad c=p_n,
 \label{eq:pfe}
\end{equation}
であり、実際の分母根 \(\rho_i\) に対して
\begin{equation}
  \beta_i=-\rho_i,
  \qquad
  A_i=c\,
  \frac{\prod_{j=1}^{n}(\rho_i-r_j)}
       {\prod_{j\ne i}(\rho_i-\rho_j)}
  \label{eq:residue}
\end{equation}
となる。実装は係数配列として積を展開し、分子と分母の最高次項を相殺して
proper fraction にした後、各 pole で評価して留数を得る。
最後に shift \(\beta_i\) が小さい順へ並べ替える。

逆数については根と pole を交換して
\begin{equation}
 \frac{1}{R_{n,n}(x)}=\frac{1}{c}
 +\sum_{i=1}^{n}\frac{\widetilde A_i}{x+\widetilde\beta_i}
 \end{equation}
を同じルーチンで作る。この形は、格子 QCD で shift の異なる線形系を
multi-shift solver に渡す用途を意識したものと考えられる。

\section{コマンドラインプログラムの入出力}

実行形式は README では
\begin{verbatim}
./test y z n d lambda_low lambda_high precision
\end{verbatim}
である。すなわち
\begin{equation}
 f(x)=x^{y/z},\quad x\in[\lambda_{\rm low},\lambda_{\rm high}],
\end{equation}
を次数 \((n,d)\)、指定十進桁数相当の演算で近似する。README の例は
\begin{verbatim}
./test 1 2 5 5 0.0004 64 40
\end{verbatim}
であり、\(\sqrt{x}\) の \((5,5)\) 有理近似を作る。

\path{approx.dat} にはまず \(x^{y/z}\)、次に \(x^{-y/z}\) の部分分数係数
\(c,A_i,\beta_i\) を倍精度形式で書く。多倍長で内部計算していても、公開 API と
このファイル出力時には \code{double} へ変換されるため、出力精度はほぼ倍精度までである。

\path{error.dat} には
\begin{equation}
  x,\quad \frac{R(x)-f(x)}{f(x)}
\end{equation}
を、\(x\leftarrow1.01x\) という対数的な間隔で記録する。
ループ条件が \(x<b\) なので上端 \(b\) 自体は出力されない。

\path{main.C} は \code{force\_...dat} と \code{energy\_...dat} という名前も作るが、
その後は使用しない。

\section{多倍長演算の実装}

\code{bigfloat} は GMP の \code{mpf\_t} を C++ 演算子で包んだ薄いクラスである。
加減乗除・平方根・整数乗は GMP を使い、一般のべきと指数関数には MPFR の関数を使う。
ヘッダは古い \path{mpf2mpfr.h} 互換層に依存しており、現在の GMP/MPFR 環境へ
そのまま持ち込めるとは限らない。

指定された十進桁数 \(D\) は
\begin{equation}
  B=\left\lfloor 3.321928094D\right\rfloor
\end{equation}
bit に変換して GMP の default precision に設定される。guard bits は加えない。
また、区間端はコンストラクタへ \code{double} で渡され、\(\pi\)、三角関数、
初期点変換にも倍精度定数・関数が混在する。したがって「40 桁を指定した」場合でも、
すべての入力と前処理が 40 桁で行われるわけではない。

\section{Julia 移植前に把握すべき制約と不具合}

以下は単なる文体の古さではなく、再実装時に挙動を保存するか改善するかを
明示的に決める必要がある点である。

\begin{longtable}{@{}p{0.27\linewidth}p{0.67\linewidth}@{}}
\toprule
分類 & コード上の性質・問題 \\
\midrule
\endhead
収束判定 & spread の閾値は常に \(10^{-15}\)。要求精度に連動せず、外側反復の上限もない。\\
極値探索 & 最大 10 回の固定ステップ山登りであり、極値を高精度に bracket・精密化しない。\\
交互性 & 誤差の符号を計算するが使用せず、隣接極値が逆符号か検査しない。\\
線形代数 & 単項式基底の独自 Gaussian 消去であり、高次数や広い区間で悪条件化しやすい。\\
根探索 & 実根のみ。bracket なしの Newton 法と逐次デフレーションで、複素根・重根に弱い。\\
部分分数 & \(n=d\) のみ。重 pole や複素共役 pole を扱う設計ではない。\\
精度境界 & 入力・公開評価 API・部分分数出力が \code{double} で、多倍長結果を倍精度へ丸める。\\
入力検査 & 引数個数、\(q\ne0\)、\(0<a<b\)、非負次数などを検査しない。\\
エラー処理 & ライブラリ内部から \code{exit(0)} する箇所があり、失敗でも成功終了 status になる。\\
戻り値 & \code{getPFE}/\code{getIPFE} は失敗時にも成功時にも 0 を返し、状態を区別できない。\\
状態管理 & 再生成前に \code{foundRoots} を false に戻さないため、2 回目の根探索失敗後に古い成功状態が残りうる。\\
メモリ & \path{main.C} は \code{new[]} した配列へ \code{delete} を使い、未定義動作になる。根探索の早期 return にも一時配列の leak がある。\\
配列確保 & \code{IPS} を \code{new int[neq*sizeof(int)]} としており、必要量より \code{sizeof(int)} 倍多く確保する。\\
次数不一致 & 近似生成自体は \(n\ne d\) に対応するが、main は部分分数失敗を確認せず配列を出力する。特に \(n>d\) では範囲外アクセスの危険がある。\\
指数因子 & 式~\eqref{eq:exp-target} の \(k_j\) は符号付き整数だが、整数べき関数の実体は unsigned exponent 用であり、負の \(k_j\) は安全でない。\\
古い依存関係 & Makefile は \path{$HOME/gmp} と \path{mpf2mpfr.h} を前提にし、現代的な標準配置や package 管理に対応しない。\\
\bottomrule
\end{longtable}

さらに、main の \(x\leftarrow1.01x\) という出力ループは \(x=0\) なら進まない。
分数べきと相対誤差の定義も考えると、現在の用途では区間下端を正に制限するのが自然である。

\section{Julia 版へ移す際の推奨分割}

元コードとの数値比較を可能にするには、最初からアルゴリズムを全面的に変更するより、
まず「互換経路」と「改善経路」を分けるのが安全である。

\begin{enumerate}
  \item \textbf{対象関数と誤差}: \(f\)、相対誤差、区間を独立した値・関数として定義する。
  \item \textbf{係数生成}: 零点から式~\eqref{eq:linear-row} を組み立てる純粋関数にする。
  \item \textbf{交換反復}: 極値位置、零点、spread、反復履歴を結果として保存し、
        最大反復数と要求精度に応じた停止条件を持たせる。
  \item \textbf{後処理}: 根探索と部分分数展開を係数生成から分離する。
        近似係数が必要なだけなら、根探索失敗で全計算を失敗させない。
  \item \textbf{精度}: \code{setprecision(BigFloat, bits) do ... end} の局所 scope を使い、
        区間端も文字列から \code{BigFloat} として parse して精度を保つ。
  \item \textbf{結果型}: 係数、誤差、収束判定、極値列、根、留数を named result に格納し、
        \code{Float64} 変換は利用者が明示した場合だけ行う。
\end{enumerate}

互換性確認では、少なくとも次を自動テストすべきである。

\begin{itemize}
  \item 各補間点で \(|R(\xi_i)-f(\xi_i)|\) が作業精度に応じて十分小さいこと。
  \item 区間内に pole がないこと。
  \item \(K+1\) 個の極値で符号が交互になり、絶対値の spread が許容値以下であること。
  \item 係数形式、因数形式、部分分数形式が同じ値を返すこと。
  \item \(R(x)\) と逆数近似 \(1/R(x)\) が相互に整合すること。
  \item README の \((p,q,n,d,a,b)=(1,2,5,5,0.0004,64)\) を回帰例にすること。
\end{itemize}

改善版では、極値探索を bracket 付き一次元最適化に替え、誤差符号の交互性を
収束条件へ含めることが最優先である。その次に、区間の正規化と安定な多項式基底、
複素根を扱える固有値ベースの根計算、重 pole を含む部分分数処理を検討できる。

\section{コード位置との対応}

主要な説明と元コードの対応は次のとおりである（行番号は調査時点）。

\begin{center}
\begin{tabular}{@{}lll@{}}
\toprule
処理 & ファイル & 行 \\
\midrule
反復全体 & \path{upstream/src/alg_remez.C} & 110--179 \\
初期点 & \path{upstream/src/alg_remez.C} & 264--299 \\
極値探索と零点更新 & \path{upstream/src/alg_remez.C} & 301--383 \\
補間方程式 & \path{upstream/src/alg_remez.C} & 385--423 \\
Horner 評価と相対誤差 & \path{upstream/src/alg_remez.C} & 425--473 \\
Gaussian 消去 & \path{upstream/src/alg_remez.C} & 475--588 \\
根探索 & \path{upstream/src/alg_remez.C} & 590--682 \\
部分分数展開 & \path{upstream/src/alg_remez.C} & 684--752 \\
コマンドラインと出力 & \path{upstream/src/main.C} & 13--101 \\
多倍長 wrapper & \path{upstream/src/bigfloat.h} & 16--209 \\
\bottomrule
\end{tabular}
\end{center}

\section{要約}

この実装の中心は、\(K=n+d+1\) 個の誤差零点で \(P/Q=f\) を解く線形補間と、
零点で区切られた \(K+1\) 区間の相対誤差極値を等振幅にする零点移動である。
収束後の根・部分分数展開まで含め、格子 QCD で利用しやすい
\(c+\sum_i A_i/(x+\beta_i)\) を生成する。

Julia 化では、まずこの零点交換法を再現して既存結果との回帰比較を可能にし、
その上で極値探索、収束検証、線形代数、根探索を段階的に置き換えるのがよい。
特に「relative minimax」「分母 monic」「零点補間」「\(n=d\) の部分分数」という
四つが元実装の数値仕様の核である。

\end{document}
